\section{Problem formulation}\label{sec:problem}

\paragraph{Swarm and artifacts.}
A swarm is a set $\mathcal{A}$ of agents that act over time on a set of shared \emph{artifacts} (wiki pages, notes, memory entries, chat rooms) and produce \emph{outputs} (a value, a link, a patch). An artifact has versions; a version is a sequence of \emph{chunks}, each written by one agent. A \emph{reader} that opens an artifact is \emph{served} a copy of its current version.

\paragraph{Logs.}
The \emph{edit log} $E$ holds one record per write or output: (agent, time, channel, content), where the channel is the room, page or thread in which the record is visible. The \emph{read log} $\mathcal{R}$, when it exists, holds one record per served item: (reader, time, channel, the exact content served, the author of the chunk it came from). Most deployments keep $E$ and not $\mathcal{R}$. A \emph{registry} (Section~\ref{sec:system}) is a third, optional record kept by a gateway.

\begin{definition}[Exposure copy and source]\label{def:copy}
Agent $a$ \emph{copies} an item $u$ from agent $b$ if $a$ was served a copy of $u$ whose chunk was written by $b$ before $a$ first used $u$. The \emph{immediate source} of $a$'s use of $u$ is the author of the served copy it used. An agent that used $u$ without having been served it \emph{worked alone}.
\end{definition}
This is the standard exposure notion in contagion studies: it says nothing about intent, and it is what an operator can act on (the agent saw the item). Where several served copies are byte-identical, the immediate source is not determined by the text alone; we return to this in Section~\ref{sec:setup}.

\begin{definition}[Tracer and scores]\label{def:tracer}
A \emph{tracer} maps the logs it is allowed to see to a prediction for each agent: an earlier agent as source, or ``independent''. On a set of agents with known truth, the \emph{top-1 accuracy} is the share of copying agents whose source is named correctly; \emph{recall} is the share of copying agents accused of copying at all; the \emph{false-accusation rate} is the share of agents that worked alone but were blamed on someone. The \emph{predicted spread graph} is the graph of predicted sources.
\end{definition}

\begin{definition}[Carrier coverage]\label{def:coverage}
For a copy event, its \emph{carrier} is an earlier user of the same distinctive item that is visible in the log in the channel where the copy happens. \emph{Carrier coverage} is the share of copy events that have at least one visible carrier. It measures how much of the propagation channel the log shows, and is the audit's headline.
\end{definition}

\begin{definition}[Re-check set]\label{def:recheck}
Given a bad item entering through agent (or post) $o$, the \emph{affected} outputs are those whose chain of copying leads back to $o$. A \emph{re-check set} is a list of outputs predicted to be affected. It is scored by its \emph{size} as a share of all outputs (the work), by \emph{recall} of the affected outputs, and by \emph{precision}.
\end{definition}

\paragraph{What we assume.}
(A1)~Logs are honest: agents do not tamper with the edit log. (A2)~Agents do not deliberately strip or forge tokens; we study what \emph{unintended} copying does to them (Section~\ref{sec:limits}). (A3)~The operator controls the gateway and the shared artifacts, but not what the agents' models do with them. We do not model adversarial agents.

\paragraph{The questions.}
(Q1)~How much of the copying in a real swarm can \emph{any} tracer recover from $E$ alone? (Q2)~How should copying be called without being fooled by shared habits? (Q3)~If an operator can change how artifacts are served, what should be changed and how well does it work across models? (Q4)~After a bad item spreads, what has to be re-checked?
